% ملف: main.tex
\documentclass[12pt]{article}

\usepackage{fontspec}        %(c) It's importent with XeLaTeX
\usepackage{polyglossia}     %(c) Languages management \textendash and \textarabic.

\setmainlanguage{arabic}
\setotherlanguage{english}
%(c) choos an arabic font has been instulled in your system: 
\newfontfamily\arabicfont[Script=Arabic,Scale=1.0]{Amiri} % And this font is good (Scheherazade New)
%(c) If you want to mix English text within Arabic text:
\newfontfamily\englishfont{Latin Modern Roman}
\newcommand{\al}{\textarabic} %(c) to make the command easier
\newcommand{\el}{\textenglish} %(c) to make the command easier

\begin{document}

\title{تجربة \LaTeX{} \textenglish{(XeLaTeX + polyglossia)} بالعربية}
\author{محمد علي قطان}
\date{\today}
\maketitle

\section{مقدمة}
% If you encounter any problems while writing in Arabic and English, you can use these definitions (\textarabic{Type in arabic here}, \text{Type in English here}) to ensure that no problems or interference between the two languages ​​occur.
هذه تجربة بسيطة لاختبار دعم اللغة العربية داخل \LaTeX{} باستخدام \textenglish{\textbf{XeLaTeX}} و \textenglish{\textbf{polyglossia}}.
% يفضل أن لا تستخدم التعريفات للعربي والنكليزي معا، اللغة الأساسية أكتبها بدون تعريفات واللغة الدخيلة اكتبها بتعرفات هذه افضل استخدام

\subsection{مزج العربية والإنجليزية}
نستطيع كتابة English داخل النص العربي بسهولة: \textenglish{This is English inside Arabic text}.

\end{document}
